\documentclass[11pt]{article}
\usepackage[margin=1in]{geometry}
\usepackage{amsmath,amssymb,amsthm}
\usepackage{enumitem}
\usepackage{hyperref}

\newtheorem{theorem}{Theorem}
\newtheorem{proposition}[theorem]{Proposition}
\newtheorem{corollary}[theorem]{Corollary}
\newtheorem{remark}[theorem]{Remark}

\newcommand{\Zm}{\mathbb Z_m}
\newcommand{\Xm}{X_m}
\newcommand{\rhoc}[1]{\rho_{#1}}
\newcommand{\SelCorr}{\mathrm{SelCorr}}
\newcommand{\SelStar}{\mathrm{SelStar}}

\title{d5\_234: Graph-side split --- \(G2\) closes by cyclic transport, while \(G1\) remains a genuine splice theorem}
\author{}
\date{2026-03-21}

\begin{document}
\maketitle

\section{Purpose}
The current manuscript leaves two graph-side frontier objects:
\begin{itemize}[leftmargin=2.4em]
\item \(G1\): colorwise selector-compatibility;
\item \(G2\): remaining-color transfer.
\end{itemize}
This note separates them cleanly.
The main conclusion is that \(G2\) is no longer an independent frontier theorem.
It follows formally from cyclic color-equivariance once a \(G1\)-compatible colorwise package exists.
By contrast, \(G1\) remains genuine: the surfaced mixed witness already gives
colorwise permutations, but its naive cyclic package fails pointwise outgoing
exhaustiveness at an explicit vertex for every odd \(m>2\).

\section{Abstract cyclic transport theorem}
Fix
\[
\Xm := (\Zm)^5.
\]
For \(c\in\mathbb Z/5\mathbb Z\), let \(\rhoc{c}:\Xm\to\Xm\) be the cyclic
coordinate shift
\[
\bigl(\rhoc{c}(x)\bigr)_i := x_{i+c}
\qquad (i\in\mathbb Z/5\mathbb Z).
\]
Suppose a color-relative local rule is given by:
\begin{enumerate}[label=\textup{(\arabic*)},leftmargin=2.8em]
\item a local feature map \(F\) on the color-$0$ chart;
\item a relative direction selector
      \(a\colon \mathrm{Im}(F)\to\{0,1,2,3,4\}\);
\item the color-$c$ absolute direction rule
      \[
      d_c(x) := a\bigl(F(\rhoc{c}x)\bigr)+c \pmod 5.
      \]
\end{enumerate}
Define
\[
 g_c(x):=x+e_{d_c(x)}.
\]

\begin{proposition}[cyclic conjugacy of color-relative rules]
\label{prop:cyclic-conjugacy}
For every color \(c\),
\[
 g_c = \rho_{-c}\circ g_0\circ \rho_c.
\]
\end{proposition}

\begin{proof}
Fix \(x\in X_m\) and put \(y:=\rho_c x\).
By definition,
\[
 g_0(y)=y+e_{a(F(y))}.
\]
Applying \(\rho_{-c}\) gives
\[
 \rho_{-c}g_0\rho_c(x)=x+e_{a(F(\rho_c x))+c}=x+e_{d_c(x)}=g_c(x).
\]
\end{proof}

\begin{corollary}[graph-side transfer = conjugacy]
\label{cor:G2-closed}
Assume a \(G1\)-compatible colorwise package \(g_0,\dots,g_4\) is obtained from
one color-relative local rule family.
Then every conjugacy-invariant property of the closed color-$4$ branch transfers
to every other color.
In particular, permutation, cycle-structure, first-return, grouped-return,
nested-return, and Hamilton-type properties transport from color $4$ to colors
\(0,1,2,3\).
Thus the manuscript object \(G2\) is not an additional frontier theorem; it is
an immediate corollary of \(G1\) plus cyclic color-equivariance.
\end{corollary}

\begin{proof}
By Proposition~\ref{prop:cyclic-conjugacy}, each \(g_c\) is conjugate to \(g_0\),
and therefore also to \(g_4\).
Every listed property is invariant under conjugation of the global map and under
the corresponding transport of the chosen section/return coordinates.
\end{proof}

\section{Why \(G1\) is still genuine}
The surfaced bundle contains an exact cyclic family coming from the canonical
mixed witness \texttt{mixed\_008}, with payload
\[
(\texttt{layer2}=0/2,\ \texttt{pred\_sig1\_wu2},\ \texttt{layer3 modes }0/3,3/0).
\]
This family already satisfies colorwise incoming uniqueness, hence each color map
is a torus permutation on the checked moduli.
Nevertheless it does \emph{not} solve \(G1\), because pointwise outgoing
exhaustiveness fails.

\begin{proposition}[explicit failure of naive outgoing exhaustiveness]
\label{prop:mixed008-fails-G1}
For every odd \(m>2\), the naive cyclic package associated to the surfaced
\texttt{mixed\_008} witness fails pointwise outgoing exhaustiveness at
\[
 x=(0,0,0,1,2)\in X_m.
\]
More precisely, the five chosen absolute directions at \(x\) are
\[
[3,1,0,3,4],
\]
so direction \(2\) is missing and direction \(3\) occurs twice.
\end{proposition}

\begin{proof}
At the chosen vertex, all five relevant states lie in layer $3$.
The layer-$3$ old bit is \(q=-1\), and because the color-relative
\(q\)-coordinates at \(x\) are
\[
(0,0,1,2,0),
\]
all five old-bit values are $0$ when \(m>2\).
Hence the mixed witness uses only the first entry of the active layer-$3$ mode:
\begin{itemize}[leftmargin=2.4em]
\item if the predecessor flag \(p=0\), mode \(0/3\) chooses relative direction \(0\);
\item if the predecessor flag \(p=1\), mode \(3/0\) chooses relative direction \(3\).
\end{itemize}
For the flag \texttt{pred\_sig1\_wu2}, the predecessor in relative direction $1$
changes only the color-relative \(q\)-coordinate, so the flag value is exactly
\(1\) iff the color-relative sum \(w+u\) equals $2$.
At the chosen vertex the color-relative pairs \((w,u)\) are
\[
(0,2),\ (1,0),\ (2,0),\ (0,0),\ (0,1),
\]
so the corresponding flag values are
\[
(1,0,1,0,0).
\]
Therefore the relative directions are
\[
(3,0,3,0,0).
\]
Adding the color index \(c\) modulo $5$ yields the absolute directions
\[
(3,1,0,3,4),
\]
exactly as claimed.
These five directions do not exhaust \(\{0,1,2,3,4\}\).
\end{proof}

\begin{remark}
So the mixed witness already gives the \emph{transport} side of the graph story
(colorwise conjugacy and colorwise permutations), but not the \emph{selector}
side.
This is why \(G1\) cannot be replaced by a naive cyclic copy of the current
mixed rule.
A genuine splice/compatibility theorem is still needed.
\end{remark}

\section{Current honest graph-side bookkeeping}
After the front-end cleanup and the present split, the graph side should be read
as follows.
\begin{enumerate}[label=\textup{(\arabic*)},leftmargin=2.8em]
\item \textbf{Closed now:} \(G2\), in the precise sense of
      Corollary~\ref{cor:G2-closed}.
\item \textbf{Still live:} \(G1\), i.e.
      the selector-compatibility/splice theorem that inserts the already closed
      color-$4$ branch into a truly outgoing-exhaustive colorwise package.
\item \textbf{Missing concrete data for full closure of \(G1\):}
      an explicit surfaced model of the corrected selector baseline
      \(\SelCorr\), together with the local defect certificate showing how the
      color-$4$ branch \(\SelStar\) is spliced without breaking pointwise
      outgoing exhaustiveness or colorwise incoming uniqueness.
\end{enumerate}
In particular, the open graph-side list should now shrink from \(G1,G2\) to
just \(G1\).

\end{document}
